% !TEX root = alder2021leakage.tex
% FICTIONAL TEST MATERIAL: this paper does not exist; it was written for the
% citation-audit fixture.
\documentclass{article}
\usepackage{amsmath}

\title{Acoustic Leak Localization in Cast-Iron Water Mains}
\author{A. Fixture \and B. Fixture}
\date{Journal of Fixture Studies, 2021}

\begin{document}
\maketitle

\noindent\textbf{FICTIONAL TEST MATERIAL: this paper does not exist; it was
written for the citation-audit fixture.} No author, utility, venue, or
measurement described here is real.

\begin{abstract}
Leaks on buried water mains are expensive to find because excavation is slow
and a wrong dig is costly. We evaluate acoustic leak-noise correlators on a
fictional municipal network and report how accurately they locate leaks on
cast-iron mains up to 300~mm in diameter. Across 140 confirmed leaks, the
correlators located 92\% of them, with a median localization error of 1.5~m.
We also compare the localization error on cast-iron, ductile-iron, and PVC
mains, and find that the error grows as the pipe wall becomes less rigid.
\end{abstract}

\section{Introduction}
\label{sec:introduction}

A leak-noise correlator places two acoustic sensors on fittings on either side
of a suspected leak. The noise of escaping water travels along the pipe wall
and through the water column to both sensors. The difference in arrival times,
together with the distance between the sensors and the speed of sound in the
pipe, gives the position of the leak between them. The method is old, but its
accuracy depends strongly on the pipe material and diameter, and published
field figures are scattered.

This paper reports a field evaluation on the fictional Fixture Valley network.
Our question is narrow: how often, and how precisely, do correlators locate
leaks on cast-iron mains, and how does that compare with other common pipe
materials?

\section{Method}
\label{sec:method}

\subsection{Network and leak set}
\label{sec:method-network}

The evaluation covered cast-iron distribution mains with nominal diameters from
100~mm to 300~mm. A leak entered the study when a survey crew suspected it and
a later excavation confirmed it. In total 140 leaks on cast-iron mains were
confirmed during the study period. A smaller comparison set of leaks on
ductile-iron and PVC mains was collected in the same way.

\subsection{Correlation procedure}
\label{sec:method-procedure}

For each suspected leak, the crew placed two sensors on hydrants or valves no
more than 200~m apart. The correlator computed the cross-correlation of the two
signals and reported the leak position. The crew then marked that position on
the ground, and the excavation team recorded the true leak position after
excavation. The localization error is the distance between the marked and the
true position. A leak counts as \emph{located} when the excavation at the
marked position exposed it, which we define as an error of 3~m or less.

\section{Results}
\label{sec:results}

\subsection{Localization rate and error on cast iron}
\label{sec:results-castiron}

The acoustic correlators located 92\% of the 140 confirmed leaks on cast-iron
mains up to 300~mm in diameter. The median localization error on these mains
was 1.5~m. Errors were smallest on the 100~mm and 150~mm mains and grew
slightly with diameter, but the median stayed below 2~m in every diameter
class. The leaks the correlators failed to locate were mostly at joints close
to a branch, where the noise reached the sensors by more than one path.

\subsection{Localization error by pipe material}
\label{sec:results-material}

Table~\ref{tab:material} gives the median localization error by pipe material.
Cast iron gave the smallest error, at 1.5~m. Ductile iron followed, at 2.1~m.
PVC gave by far the largest error, at 4.8~m, because its flexible wall damps
the leak noise and lowers the speed of sound along the pipe, so a small timing
error turns into a large position error.

\begin{table}[h]
  \centering
  \caption{Median leak-localization error by pipe material
    (Section~\ref{sec:results-material}).}
  \label{tab:material}
  \begin{tabular}{lr}
    \hline
    Pipe material & Median localization error \\
    \hline
    Cast iron     & 1.5 m \\
    Ductile iron  & 2.1 m \\
    PVC           & 4.8 m \\
    \hline
  \end{tabular}
\end{table}

\subsection{Practical consequence}
\label{sec:results-practice}

On cast iron, a median error of 1.5~m means that a single standard excavation
usually exposes the leak. On PVC, the crew should expect to widen the dig or to
confirm the position with a second sensor spacing.

\section{Limitations}
\label{sec:limitations}

This study measures leak localization only. It has clear limits on what it
can support.

This study does not examine pressure management: network pressure was not
changed during the study, and we draw no conclusion about how pressure affects
leaks. This study does not examine bursts: every event in the data set was a
background or reported leak, not a main break, and burst frequency was not
recorded. This study does not examine water quality: no samples were taken and
no water-quality parameter was measured. Finally, the comparison sets for
ductile iron and PVC are smaller than the cast-iron set, so their medians are
less certain.

\end{document}
